\documentclass[11pt,a4paper]{ctexart}
\usepackage[margin=2.25cm]{geometry}
\usepackage{amsmath,amssymb,amsthm,mathtools,booktabs,array,longtable,hyperref,microtype}
\usepackage{enumitem}
\hypersetup{colorlinks=true,linkcolor=black,urlcolor=blue,citecolor=black}
\setlength{\parindent}{2em}
\setlength{\parskip}{0.25em}
\linespread{1.08}
\newtheorem{theorem}{定理}
\newtheorem{lemma}{引理}
\newtheorem{corollary}{推论}
\renewcommand{\proofname}{证明}
\newcommand{\vthree}{v_3}
\newcommand{\Fthree}{F_3}
\newcommand{\Qk}{\mathcal Q_k}
\newcommand{\Cfrak}{\mathfrak C}
\title{\bfseries 固定伸缩素数族中的\\高阶 $F_3$ 相互作用}
\author{研究札记 · F3-HI-1 / F3-HI-2 / F3-HI-3}
\date{2026 年 9 月 29 日}
\begin{document}
\maketitle

\begin{abstract}
设 $\Fthree(m)=\vthree(m+1)-\vthree(m-1)$。本文研究固定伸缩形状
\[
 p_0=n,\qquad p_1=n+2\cdot 3^{k+1}d,\qquad p_2=n+6\cdot 3^{k+1}d.
\]
在要求 $d,p_0,p_1,p_2$ 全为素数且 $\Fthree(n)=k\ge1$ 时，三个单点的 $F_3$ 值恒为 $k$，三个两两乘积的 $F_3$ 值恒为 $-k$，但三因子乘积的 $F_3$ 值仍可达到任意预先指定的层级 $R\ge k+1$。更进一步，三因子深度具有显式极限分布：最低层概率为 $1/2$，其后各层依次为 $1/3,1/9,1/27,\ldots$。证明先把全部未记录的高阶信息压缩到一个二次多项式余因子 $G_n(d)$，再利用其模 $3$ 的唯一简单根做逐层提升，最后用 Green--Tao 的素数线性形式定理把有限剩余类比例转移到真实全素数配置。文末还给出一般有限剩余模式的素数伸缩实现定理。
\end{abstract}

\noindent\textbf{关键词：} 三进赋值；$F_3$；高阶相互作用；Hensel 型提升；素数线性形式；Green--Tao。

\section{定义与主定理}
对正整数 $a$，$\vthree(a)$ 表示 $3$ 在 $a$ 中的指数。对整数 $m>1$，定义
\[
\Fthree(m)=\vthree(m+1)-\vthree(m-1)\in\mathbb Z.
\]
固定 $k\ge1$，令
\[
M=3^{k+1},\qquad p_0=n,\quad p_1=n+2Md,\quad p_2=n+6Md.
\]
定义
\[
\Qk(X)=\{(n,d):X<n,d\le2X,\ d,p_0,p_1,p_2\text{ 均为素数},\ \Fthree(n)=k\}.
\]
记
\[
D(n,d)=\Fthree(p_0p_1p_2),
\quad
\Cfrak=\prod_{\substack{\ell\ge5\\ \ell\ \mathrm{prime}}}
\frac{\ell^2(\ell-3)}{(\ell-1)^3},
\quad
\kappa_k=3^{2-k}\Cfrak.
\]

\begin{theorem}[F3-HI-1]
在 $\Qk(X)$ 中恒有
\[
\Fthree(p_0)=\Fthree(p_1)=\Fthree(p_2)=k,
\]
\[
\Fthree(p_0p_1)=\Fthree(p_0p_2)=\Fthree(p_1p_2)=-k.
\]
并且
\[
|\Qk(X)|=(\kappa_k+o_k(1))\frac{X^2}{(\log X)^4}.
\]
对每个固定整数 $R\ge k+1$，
\[
\#\{(n,d)\in\Qk(X):D(n,d)=R\}
=\bigl(\kappa_k w_k(R)+o_{k,R}(1)\bigr)\frac{X^2}{(\log X)^4},
\]
其中
\[
w_k(R)=
\begin{cases}
\frac12,&R=k+1,\\[2mm]
3^{-(R-k-1)},&R\ge k+2.
\end{cases}
\]
因此每个 $R\ge k+1$ 都有无穷多个全素数实现，且偏移系数 $2M,6M$ 不随 $R$ 改变。
\end{theorem}

\begin{corollary}[F3-HI-2]
即使同时知道同一固定伸缩形状中的三个单点 $F_3$ 值、三个两两乘积 $F_3$ 值，并且步长 $d$ 也为素数，三因子乘积的 $F_3$ 值仍不能由这些二阶数据给出有限上界。
\end{corollary}

当 $k=1$ 时，$M=9$，形状成为
\[
(n,n+18d,n+54d).
\]
六个低阶值固定为 $(1,1,1)$ 与 $(-1,-1,-1)$，而三因子深度可取任意 $R\ge2$。

\section{六个低阶值为何完全锁定}
假设 $n>1$、$d>0$、$3\nmid d$ 且 $\Fthree(n)=k>0$。于是
\[
\vthree(n+1)=k,\qquad \vthree(n-1)=0,
\]
所以可写成
\[
n=-1+3^k a,\qquad 3\nmid a. \tag{2.1}
\]

\begin{lemma}
三个单点均满足 $\Fthree(p_i)=k$。
\end{lemma}
\begin{proof}
三个 $p_i$ 都满足 $p_i\equiv n\pmod{3^{k+1}}$。因此
\[
p_i+1\equiv3^k a\pmod{3^{k+1}},
\qquad p_i-1\equiv-2\pmod3.
\]
故 $\vthree(p_i+1)=k$ 且 $\vthree(p_i-1)=0$。
\end{proof}

\begin{lemma}
任意 $i\ne j$ 都有 $\Fthree(p_ip_j)=-k$。
\end{lemma}
\begin{proof}
由 $p_ip_j\equiv n^2\pmod{3^{k+1}}$ 与 (2.1)，
\[
n^2-1=-2\cdot3^ka+3^{2k}a^2
\equiv-2\cdot3^ka\pmod{3^{k+1}}.
\]
因为 $3\nmid2a$，有 $\vthree(p_ip_j-1)=k$。另一方面 $p_ip_j\equiv1\pmod3$，故 $\vthree(p_ip_j+1)=0$。于是 $\Fthree(p_ip_j)=-k$。
\end{proof}

到这里，六个低阶数据已经固定，而且完全没有使用素数性。

\section{三因子乘积中的可提升根}
定义
\[
A(n)=\frac{n^3+1}{M}.
\]
由 (2.1) 展开得
\[
A(n)=a-3^ka^2+3^{2k-1}a^3\in\mathbb Z,
\qquad A(n)\equiv a\pmod3. \tag{3.1}
\]
另一方面
\[
\begin{aligned}
p_0p_1p_2+1
&=n(n+2Md)(n+6Md)+1\\
&=n^3+1+8Mn^2d+12M^2nd^2\\
&=M G_n(d),
\end{aligned}
\]
其中
\[
\boxed{G_n(d)=A(n)+8n^2d+12Mnd^2.} \tag{3.2}
\]
三个因子都模 $3$ 同余于 $-1$，所以 $p_0p_1p_2\equiv-1\pmod3$，从而
\[
\boxed{D(n,d)=k+1+\vthree(G_n(d)).} \tag{3.3}
\]
真正没有被六个低阶值记录的信息，恰好落在 $G_n(d)$ 中。

模 $3$ 看，
\[
G_n(d)\equiv a+2d\pmod3,
\qquad
G_n'(d)=8n^2+24Mnd\equiv2\pmod3. \tag{3.4}
\]
因此模 $3$ 恰有一个根 $d\equiv a$，而且导数是单位。

\begin{lemma}[唯一根逐层提升]
对每个 $t\ge1$，$G_n(d)\equiv0\pmod{3^t}$ 在单位剩余类中恰有一个根。
\end{lemma}
\begin{proof}
若 $u$ 是模 $3^t$ 的根，则
\[
G_n(u+3^tb)
\equiv G_n(u)+3^tbG_n'(u)\pmod{3^{t+1}}.
\]
除以 $3^t$ 后，因 $G_n'(u)\not\equiv0\pmod3$，$b\bmod3$ 中恰有一个选择继续成为根。由模 $3$ 唯一根归纳即可。
\end{proof}

\begin{lemma}[精确有限比例]
在模 $3^T$ 的单位类中均匀计数，则
\[
\Pr(\vthree(G_n(d))\ge t)=\frac{1}{2\cdot3^{t-1}}
\qquad(1\le t\le T).
\]
因此
\[
\Pr(\vthree(G_n(d))=0)=\frac12,
\qquad
\Pr(\vthree(G_n(d))=t)=3^{-t}\quad(t\ge1).
\]
\end{lemma}
\begin{proof}
模 $3^T$ 一共有 $2\cdot3^{T-1}$ 个单位类。唯一的模 $3^t$ 根有 $3^{T-t}$ 个提升，因此门槛比例为
\[
\frac{3^{T-t}}{2\cdot3^{T-1}}=\frac1{2\cdot3^{t-1}}.
\]
相邻门槛相减得到精确层比例。
\end{proof}

结合 (3.3)，局部深度分布正是
\[
\Pr(D=k+1)=\frac12,
\qquad
\Pr(D=k+1+t)=3^{-t}\quad(t\ge1). \tag{3.5}
\]
这只是有限剩余类的精确比例，并没有把素数事件当作独立随机事件。

\section{四素数系统与奇异乘积}
固定 $A\ge k+1$，令 $Q=3^A$，并固定一个允许的模 $Q$ 参数类。写
\[
n=b+Qx,\qquad d=c+Qy.
\]
需要同时为素数的四个形式为
\[
d,\quad n,\quad n+2Md,\quad n+6Md.
\]
它们的齐次方向是
\[
(0,1),\quad(1,0),\quad(1,2M),\quad(1,6M),
\]
两两不成比例。因此按 Green--Tao 的复杂度定义，该系统复杂度至多为 $2$。

局部因子可直接计算。模 $3$ 时四形式均为单位，所以
\[
\beta_3=(3/2)^4.
\]
模 $2$ 时要求 $n,d$ 都为奇数，故
\[
\beta_2=4.
\]
若 $\ell\ge5$ 为素数，要求 $d\ne0$ 后，$n/d$ 只需避开
\[
0,\quad-2M,\quad-6M.
\]
三类互异，因此允许参数数为 $(\ell-1)(\ell-3)$，于是
\[
\beta_\ell=
\frac{\ell^2(\ell-3)}{(\ell-1)^3}
=1-\frac{3\ell-1}{(\ell-1)^3}
=1+O(\ell^{-2}). \tag{4.1}
\]
所以
\[
\Cfrak=\prod_{\ell\ge5}\beta_\ell>0,
\qquad
\mathfrak S=\beta_2\beta_3\Cfrak=\frac{81}{4}\Cfrak>0. \tag{4.2}
\]
这里严格使用了 $\sum|\beta_\ell-1|<\infty$，而不是简单地用“每个局部因子为正”来推出无穷乘积为正。

\section{从有限参数类转移到真实全素数配置}
在原盒 $X<n,d\le2X$ 内，一个固定模 $Q$ 的参数类对应面积
\[
\frac{X^2}{Q^2}+O(X)
\]
的二维矩形。由 Green--Tao 关于复杂度至多 $2$ 的素数线性形式无条件定理，每个允许参数类的全素数计数为
\[
\left(\frac{\mathfrak S}{Q^2}+o_{k,Q}(1)\right)
\frac{X^2}{(\log X)^4}. \tag{5.1}
\]
主系数不依赖于具体允许类。

模 $Q=3^A$ 下，满足 $\vthree(n+1)=k$ 的 $n$ 类数为
\[
2\cdot3^{A-k-1},
\]
单位 $d$ 类数为
\[
2\cdot3^{A-1}.
\]
故允许参数类总数为
\[
4\cdot3^{2A-k-2}. \tag{5.2}
\]
对 (5.1) 求和，得到
\[
\begin{aligned}
|\Qk(X)|
&=\left(
4\cdot3^{2A-k-2}\cdot3^{-2A}\cdot\frac{81}{4}\Cfrak+o_k(1)
\right)
\frac{X^2}{(\log X)^4}\\
&=(3^{2-k}\Cfrak+o_k(1))\frac{X^2}{(\log X)^4}.
\end{aligned}
\]
即主项常数正是 $\kappa_k=3^{2-k}\Cfrak$。

固定 $R\ge k+1$，取 $A\ge R+1$。此时模 $3^A$ 信息足够判断 $\vthree(G_n(d))=R-k-1$。第 3 节已经说明，在每个允许 $n$ 类上，对应的 $d$ 类比例都相同，所以只需把精确有限比例乘进 (5.1)。得到
\[
\#\{(n,d)\in\Qk(X):D=R\}
=
(\kappa_kw_k(R)+o_{k,R}(1))\frac{X^2}{(\log X)^4}.
\]
每个 $w_k(R)>0$，因此每个指定深度都有无穷多个全素数实现。定理 F3-HI-1 与推论 F3-HI-2 得证。

\section{一般有限模式的伸缩实现}
\begin{theorem}[F3-HI-3]
固定 $s\ge2$、$A\ge1$，令 $Q=3^A$。给定模 $Q$ 的单位类
\[
a_0,\ldots,a_{s-1},b,
\]
存在固定整数
\[
0=h_0<h_1<\cdots<h_{s-1}
\]
和常数 $c>0$，使满足
\[
(n,d)\equiv(a_0,b)\pmod Q,\qquad X<n,d\le2X,
\]
且
\[
d,n+h_0d,\ldots,n+h_{s-1}d
\]
全部为素数的参数对数为
\[
(c+o(1))\frac{X^2}{(\log X)^{s+1}}.
\]
并且 $n+h_i d\equiv a_i\pmod Q$。
\end{theorem}

\begin{proof}
令
\[
P=\prod_{\substack{\ell\le s\\ \ell\text{ prime},\ \ell\ne3}}\ell.
\]
对 $i\ge1$，由中国剩余定理选取 $h_i$ 满足
\[
h_i\equiv(a_i-a_0)b^{-1}\pmod Q,
\qquad h_i\equiv0\pmod P,
\]
再分别加上足够大的 $QP$ 倍数使其严格递增。这样保证端点落入指定模 $Q$ 类。

系统 $d,n+h_0d,\ldots,n+h_{s-1}d$ 的齐次方向两两不平行，故复杂度有限。模 $3$ 时所有形式是指定单位。对 $\ell\ne3$，要求 $d\ne0$ 后，$n/d$ 需避开 $-h_i$。若 $\ell\le s$，由于 $h_i\equiv0\pmod\ell$，只排除一个类；若 $\ell>s$，最多排除 $s<\ell$ 个类。因此所有局部因子为正。除有限个整除某个 $h_i-h_j$ 的素数外，$h_i$ 两两不同，尾部局部因子为
\[
\frac{\ell^{s-1}(\ell-s)}{(\ell-1)^s}=1+O_s(\ell^{-2}),
\]
故奇异乘积正。应用一般有限复杂度的素数线性形式定理即得结论。
\end{proof}

若有限多个非空单项式
\[
Y_\alpha(x)=\prod_i x_i^{e_{\alpha,i}}
\]
在一组正整数 $x_i$ 上实现了非零目标
\[
\Fthree(Y_\alpha(x))=f_\alpha,
\]
则取 $A>\max_\alpha|f_\alpha|$，并令 $a_i\equiv x_i\pmod{3^A}$。任何落入相同模 $3^A$ 类的端点都会保持这些精确赋值，因此 F3-HI-3 给出无穷多全素数实现。

\section{一个深度 14 的有限实例}
取
\[
n=129539,\qquad d=144481.
\]
则
\[
p_1=2730197,\qquad p_2=7931513,
\]
并且四个数 $d,n,p_1,p_2$ 均为素数。三个单点值全为 $1$，三个两两乘积值全为 $-1$，而
\[
129539\cdot2730197\cdot7931513+1
=3^{14}\cdot586479720520,
\]
最后余因子不被 $3$ 整除，所以
\[
\Fthree(129539\cdot2730197\cdot7931513)=14.
\]
该实例仅用于 sanity check；无限结论完全来自前述证明。

\section{边界：为什么这不是固定间距结果}
若固定 $d=D$，则
\[
n,\qquad n+2MD,\qquad n+6MD
\]
变成单变量仿射形式，其齐次部分互相成比例，已经不属于上述有限复杂度非退化系统。因此本文不能推出固定 $D$ 时有无穷多解，也不推出孪生素数猜想或任何固定间距素数簇结论。

本文证明的是：在同一个固定\emph{形状}中，允许素数步长 $d$ 变化时，低阶 $F_3$ 数据完全相同而高阶乘法深度仍可无界变化，并且这种高阶自由度具有明确的渐近统计规律。

\section*{文稿状态与来源}
本结果在仓库中标记为 \texttt{PAPER-AUDITED}。本文是人类可读的书面证明；尚未形式化为 Lean 主定理。Markdown 证明包包含精确定理陈述、完整证明、引理依赖与审计记录、形式化映射及阻塞项。PDF 仅是稳定证明的排版版本，不替代 Lean kernel 核验。

\section*{参考文献}
\begin{thebibliography}{9}
\bibitem{GT10}
Ben Green and Terence Tao.
\emph{Linear equations in primes}.
Annals of Mathematics, 171 (2010), 1753--1850.
DOI: 10.4007/annals.2010.171.1753.

\bibitem{GT12}
Ben Green and Terence Tao.
\emph{The Möbius function is strongly orthogonal to nilsequences}.
Annals of Mathematics, 175 (2012), 541--566.

\bibitem{GTZ}
Ben Green, Terence Tao and Tamar Ziegler.
\emph{An inverse theorem for the Gowers $U^{s+1}[N]$-norm}.
Annals of Mathematics, 176 (2012), 1231--1372.
\end{thebibliography}

\end{document}
